\par\noindent\textbf{(1)} 单轴非寻常波的相速度满足 1/n\_eff²(θ)=cos²θ/n\_o²+sin²θ/n\_e²（θ是波法线与光轴夹角），来自 k\_⊥²/n\_e²+k\_∥²/n\_o²=(ω/c)²。故 n\_eff²=n\_o²n\_e²/(n\_e²cos²θ+n\_o²sin²θ)。
\par\noindent\textbf{(2)} 固定ω的波矢椭球上，群速度沿法向∇\_kω，因此在与光轴垂直和沿光轴的方向上分别按 n\_e\textasciicircum{}\{-2\}、n\_o\textasciicircum{}\{-2\} 加权。能流夹角 θ\_S 满足 tanθ\_S=(n\_o²/n\_e²)tanθ，因两主介电轴不同，θ\_S一般不等于θ，走离角 ρ=θ\_S−θ。
\par\noindent\textbf{(3)} 当相对双折射 δ=(n\_e−n\_o)/n\_o 很小时，n\_o²/n\_e²≈1−2δ。令 tanθ\_S≈(1−2δ)tanθ，线化 arctan 导数，得到 ρ≈−2δ tanθ/(1+tan²θ)=−δ sin2θ。其幅度在约45°附近最大，表明微弱光学各向异性可产生一阶横向能流偏移。
\par\noindent\textbf{(4)} 当 θ=0 或 π/2，ρ=0，能流与波法线沿主轴，且沿光轴两偏振退化。正晶体 n\_e>n\_o 时在该约定下ρ为负，负晶体反号；对于强双折射须使用完整arctan关系，不能让小参数公式在任意δ下外推。
